(1) 在折射率加权角度表示中，曲率折射面矩阵为 \(S(\Phi)=\begin{pmatrix}1&0\\-\Phi&1\end{pmatrix}\)。透镜内部厚度 \(d\) 的平移矩阵为 \(P(d,n)=\begin{pmatrix}1&d/n\\0&1\end{pmatrix}\)。整个左至右光线矩阵是 \(M=S(\Phi_2)P(d,n)S(\Phi_1)\)，矩阵按光线先经历的元件放在右侧相乘。三个矩阵行列式均为1，表达光线相空间的辛结构。

(2) 连乘得到 \(A=1-d\Phi_1/n,\ B=d/n,\ C=-\Phi_1-\Phi_2+d\Phi_1\Phi_2/n,\ D=1-d\Phi_2/n\)。定义总屈光力 \(\Phi_{\rm eff}=-C=\Phi_1+\Phi_2-(d/n)\Phi_1\Phi_2\)。由于入射和出射均为空气，等效焦距从主平面量起为 \(f=1/\Phi_{\rm eff}\)。必须保留厚度交叉项；简单将两个球面屈光力相加仅在 \(d\to0\) 时成立。

(3) 入射平行光的 \(n\theta_{\rm in}=0\)，第二面后再传播距离 \(s_b\) 成像时 \(y_{\rm image}=(A+s_bC)y_{\rm in}\)。令系数为零得后顶点焦距 \(s_b=-A/C=A/\Phi_{\rm eff}\)。同理从第一面前方轴上点发出的光，传播 \(s_f\) 后使透镜后出射光平行，条件 \(Cs_f+D=0\)，故前顶点距 \(s_f=D/\Phi_{\rm eff}\)（向左量正）。这两个顶点焦距通常不等于等效焦距 \(f\)。

(4) 在 \(d\to0\) 时，\(A,D\to1,\ B\to0\)，\(\Phi_{\rm eff}\to(n-1)(1/R_1-1/R_2)\)，故 \(s_f,s_b\to f\)，恢复薄透镜制镜公式。厚透镜可分解为物方主平面平移、单薄透镜屈光力和像方主平面平移，因而焦点从主平面测量才是相同的 \(f\)。例如若第一面折射力较大，\(A=1-d\Phi_1/n\) 可显著偏离1，后焦点相对于第二面发生位移。
